\documentclass[11pt]{article}
\usepackage[margin=1in]{geometry}
\usepackage{amsmath,amssymb,amsthm,mathtools}
\usepackage{graphicx}
\usepackage{microtype}
\usepackage[hidelinks]{hyperref}

\newtheorem{theorem}{Theorem}
\newtheorem{lemma}{Lemma}
\newcommand{\R}{\mathbb R}
\newcommand{\vol}{\operatorname{vol}}

\title{The facet is exactly minimal among all two-level\\
central sections of the regular simplex}
\author{Autonomous proof packet for arXiv:2403.13224}
\date{11 August 2026}

\begin{document}
\maketitle

\begin{center}
\fbox{\parbox{0.91\textwidth}{\textbf{Status.}
Likely valid substantial partial result.  The theorem below proves the
conjectured sharp minimum for every central normal having exactly two
coordinate values, in every dimension.  It does not prove the conjecture for
an arbitrary normal.}}
\end{center}

\section{Source question and known boundary}

Let
\[
 \Delta_n=\bigl\{x\in\R^{n+1}:x_j\geq0,\ \textstyle\sum_jx_j=1\bigr\}
\]
be the regular $n$-simplex of side length $\sqrt2$.  For a unit vector
$a\perp(1,\ldots,1)$, put
\[
 T(a)=\Delta_n\cap\{x:\langle a,x\rangle=0\}.
\]
Tang asks whether the central hyperplane parallel to a facet always gives the
minimum possible $(n-1)$-volume~\cite{Tang2024}.

\begin{figure}[ht]
\centering
\includegraphics[width=\textwidth]{figures/open_problem_crop.png}
\caption{Question 16 on page 17 of the source paper~\cite{Tang2024}.}
\end{figure}

Tang recalls Webb's probabilistic formula.  If $Y_1,\ldots,Y_{n+1}$ are
independent standard exponential random variables and $G_a$ is the density
of $\sum_j a_jY_j$, then
\begin{equation}\label{eq:webb}
 \vol_{n-1}T(a)=\frac{\sqrt{n+1}}{(n-1)!}\,G_a(0).
\end{equation}
For the facet normal,
\[
 a_{\rm fac}=\left(\sqrt{\frac n{n+1}},
 -\frac1{\sqrt{n(n+1)}},\ldots,
 -\frac1{\sqrt{n(n+1)}}\right),
\]
the density is
\begin{equation}\label{eq:facet-density}
 c_n:=G_{a_{\rm fac}}(0)=\left(\frac n{n+1}\right)^{n-1/2}.
\end{equation}

Dirksen proved the sharp bound on the entire chamber with exactly one
positive coordinate (and, by changing sign, exactly one negative coordinate),
and proved the full conjecture for $n\leq4$~\cite{Dirksen2016}.  The result
below is different: it compares all possible sign multiplicities when each
sign class is internally constant.

\section{Proof intuition}

A two-level central normal is forced by the zero-sum and unit-norm constraints
once the two multiplicities are fixed.  Its positive and negative exponential
sums are scaled Gamma variables.  Integrating the two Gamma densities at a
common positive value gives a closed formula $D_{k,m}$ for $G_a(0)$.

The comparison between adjacent multiplicities then collapses to a
one-variable monotonicity fact.  Namely,
\[
 \frac{D_{k+1,m-1}}{D_{k,m}}
 =\frac{A(k)}{A(m-1)},\qquad
 A(x)=\left(1+\frac1x\right)^{x+1/2}.
\]
The function $A$ is strictly decreasing.  Moving from the facet multiplicity
$k=1$ toward a balanced split therefore strictly increases the density and,
by~\eqref{eq:webb}, the section volume.  What remains open is controlling
normals whose positive or negative coordinates are not all equal.

\section{Exact theorem}

Call a central normal \emph{two-level} if its coordinates assume exactly two
distinct values.

\begin{theorem}[Complete two-level classification]\label{thm:main}
Let $n\geq2$, put $d=n+1$, and let $a\in\R^d$ be a two-level unit vector with
$\sum_ja_j=0$.  After a permutation and possibly replacing $a$ by $-a$, there
are integers $k,m\geq1$, $k+m=d$, $k\leq m$, such that
\[
 a=(\underbrace{\alpha,\ldots,\alpha}_{k},
    \underbrace{-\beta,\ldots,-\beta}_{m}),
 \qquad \alpha,\beta>0.
\]
Then
\begin{equation}\label{eq:Dkm}
 G_a(0)=D_{k,m}:=\binom{d-2}{k-1}
 \frac{k^{k-1/2}m^{m-1/2}}{d^{d-3/2}}.
\end{equation}
For fixed $d$, the values $D_{k,d-k}$ are strictly increasing for
$1\leq k\leq\lfloor d/2\rfloor$.  Consequently
\[
 G_a(0)\geq D_{1,n}=\left(\frac n{n+1}\right)^{n-1/2}=c_n,
\]
with equality if and only if one of the two coordinate values has
multiplicity one.  Equivalently, among all two-level central sections of
$\Delta_n$, the sections parallel to a facet are the unique minimizers, up to
the symmetries of the simplex.
\end{theorem}

\begin{proof}
The constraints $k\alpha=m\beta$ and
$k\alpha^2+m\beta^2=1$ give
\begin{equation}\label{eq:alphabeta}
 \alpha=\sqrt{\frac{m}{kd}},\qquad
 \beta =\sqrt{\frac{k}{md}}.
\end{equation}
Let $X\sim\operatorname{Gamma}(k,1)$ and
$Y\sim\operatorname{Gamma}(m,1)$ be independent.  Then
$\sum_j a_jY_j$ has the same law as $\alpha X-\beta Y$.  The density of the
latter variable at zero is the overlap integral of the densities of
$\alpha X$ and $\beta Y$:
\begin{align*}
 G_a(0)
 &=\int_0^\infty
 \frac{t^{k-1}e^{-t/\alpha}}{\Gamma(k)\alpha^k}
 \frac{t^{m-1}e^{-t/\beta}}{\Gamma(m)\beta^m}\,dt\\
 &=\frac{\Gamma(d-1)}{\Gamma(k)\Gamma(m)}
 \frac{\alpha^{m-1}\beta^{k-1}}{(\alpha+\beta)^{d-1}}.
\end{align*}
Using~\eqref{eq:alphabeta}, one has
\[
 \alpha=\frac{m}{\sqrt{kmd}},\qquad
 \beta=\frac{k}{\sqrt{kmd}},\qquad
 \alpha+\beta=\frac d{\sqrt{kmd}}.
\]
Substitution, together with
$\Gamma(d-1)/(\Gamma(k)\Gamma(m))=\binom{d-2}{k-1}$, yields
exactly~\eqref{eq:Dkm}.

It remains to compare multiplicities.  If $1\leq k<m$ and $k+m=d$, direct
cancellation in~\eqref{eq:Dkm} gives
\begin{align}\label{eq:ratio}
 \frac{D_{k+1,m-1}}{D_{k,m}}
 &=\left(1+\frac1k\right)^{k+1/2}
   \left(1-\frac1m\right)^{m-1/2}\\
 &=\frac{A(k)}{A(m-1)},
 \qquad
 A(x):=\left(1+\frac1x\right)^{x+1/2}.
\end{align}
We claim that $A$ is strictly decreasing on $(0,\infty)$.  Put $t=1/x$.
Then
\[
 \frac{d}{dx}\log A(x)
 =\log(1+t)-\frac{t(2+t)}{2(1+t)}.
\]
If
\[
 H(t)=\frac{t(2+t)}{2(1+t)}-\log(1+t),
\]
then $H(0)=0$ and
\[
 H'(t)=\frac{t^2}{2(1+t)^2}>0\qquad(t>0).
\]
Thus $(\log A)'<0$.  For every $1\leq k<\lfloor d/2\rfloor$, the adjacent
pair $k,k+1$ stays inside the half-range and one has $k<m-1$, hence
$A(k)>A(m-1)$ and~\eqref{eq:ratio} is strictly greater than one.  Therefore
$D_{k,d-k}$ is strictly increasing on that range.  Its minimum is $k=1$,
where~\eqref{eq:Dkm} becomes
\[
 D_{1,n}=\frac{n^{n-1/2}}{(n+1)^{n-1/2}}=c_n.
\]
The equality and geometric statements follow from~\eqref{eq:webb} and from
permutation/sign symmetry.
\end{proof}

\section{Verification and limitations}

The accompanying script independently checks the Gamma overlap integral,
formula~\eqref{eq:Dkm}, and ratio~\eqref{eq:ratio} at high precision for a
representative collection of multiplicities, and checks monotonicity through
$d=200$.  These calculations are sanity checks, not ingredients of the proof.

The theorem does not imply that the minimum inside every fixed sign chamber is
two-level.  In fact, Dirksen explicitly warns that this tempting reduction is
false in higher dimensions~\cite[Remark 3.4 and Lemma 3.6]{Dirksen2016}.
His balancing proposition makes all coordinates on one sign equal but leaves
the other sign ratios free.  Controlling those remaining ratios is the open
core of Tang's Question 16.

Several stronger routes were tested.  Numerical quadrature and global
optimization in dimensions $5$ through $11$ returned the facet normal to
numerical precision, and a pointwise comparison on Tang's positive contour
also survived randomized tests.  A sharpened Cauchy--Schwarz step using
$a\perp(1,\ldots,1)$ yields a genuine positive deficit, but converting that
deficit into the exact facet constant requires a new rational moment
inequality.  None of these numerical or unfinished observations is claimed as
a proof.

\section{Bounded novelty check and review recommendation}

The 11 August 2026 search covered the run's registry/solution/attempt indexes,
exact arXiv-id and keyword queries, Dirksen's arXiv:1509.06408, and the later
arXiv:2505.00944 and arXiv:2606.07163.  Dirksen's paper contains the closest
known analysis, including the one-sign theorem and examples involving
two-level normals, but the all-multiplicity monotonicity in Theorem
\ref{thm:main} was not located.  The two later papers concern maximum-section
or reverse-H\"older questions rather than this exact minimum.  Novelty
confidence is moderate: the result may be folklore because the Gamma
calculation is short.

Human review should focus on the Gamma scaling in~\eqref{eq:Dkm} and the
strictness range in~\eqref{eq:ratio}.  The packet should remain classified as
a partial result unless the arbitrary multi-level case is separately proved.

\begin{thebibliography}{9}
\bibitem{Tang2024}
C. Tang,
\emph{Simplex slicing: an asymptotically-sharp lower bound},
arXiv:2403.13224, 2024.

\bibitem{Dirksen2016}
H. Dirksen,
\emph{Sections of the regular simplex -- volume formulas and estimates},
arXiv:1509.06408; Math. Nachr. 290 (2017), 2567--2584.
\end{thebibliography}

\end{document}
